% !TeX root = bericht.tex
\section{Aufgabe 1}
\subsection{Plotten von zwei Sinussignalen}
\begin{figure}[H]
    \centering
    \includegraphics[width=0.9\textwidth]{lab1_1_0.png}
    \caption{x1 und x2}
\end{figure}

Es galt, zwei Sinussignale $x_1(t)$ und $x_2(t)$ in einem Diagramm grafisch darzustellen.

\begin{align}
    x_1(t) &= x_{1, 0} \cdot \sin(2 \pi f_1 t) \\
    x_2(t) &= x_{2,0} \cdot \sin\left(2 \pi f_2 t - \frac{\pi}{2}\right)
    \label{fig:x1-x2}
\end{align}

Mit $f_1 = 20 kHz$, $f_2 = 60 kHz$, $x_{1,0} = 2$ und $x_{2,0} = 4$. \\
Damit die Signale $x_1(t)$ und $x_2(t)$ richtig dargestellt werden, muss die Abtastfrequenz mindestens doppelt so hoch sein wie die höchste Frequenz der Signale sein (Nyquist-Shannon-Abtasttheorem).
Hier wurde eine Abtastfrequenz von $f = 10 MHz$ gewählt, welches deutlich höher ist als $2 \cdot 60kHz$ ist.
\begin{sourcecode}
    \caption{Matlab Code zur Parameterwahl und Zeitachse}
    \begin{minted}{matlab}
x1_0 = 2;
x2_0 = 4;
f1 = 20e3;
f2 = 60e3;

t = 0 : 0.1e-6 : 0.3e-3;    % Abtastfrequenz f=10MHz
    \end{minted}
\end{sourcecode}

\begin{sourcecode}
    \caption{Matlab Code zur Berechnung von $x_1(t)$ und $x_2(t)$}
    \begin{minted}{matlab}
x1 = x1_0 * sin(2*pi*f1*t);
x2 = x2_0 * sin(2*pi*f2*t - pi/2);
    \end{minted}
\end{sourcecode}

\begin{sourcecode}
    \caption{Matlab Code um die Signale zu plotten}
    \begin{minted}{matlab}
figure;
plot(t, x1)
hold on;
plot(t, x2)
hold off;

title('x_1(t) und x_2(t)');
xlabel('Zeit t in [s]');
ylabel('Amplitude');
legend('x_1 (20 kHz)', 'x_2 (60 kHz)');    
    \end{minted}
\end{sourcecode}


\subsection{Plotten von zwei Sinussignalen}
\begin{figure}[H]
    \centering
    \includegraphics[width=0.9\textwidth]{lab1_1_1.png}
    \caption{Summe der Signale x1 und x3}
    \label{fig:x1-x3}
\end{figure}
Hier zu sehen ist $x_1(t)$ und $x_3(t)$.
\begin{align}
    x_3(t) &= x_1(t) \cdot x_1(t) \\ 
    &= x_{1,0}^2 \cdot \sin^2(2 \pi f_1 t) \nonumber
\end{align}
Da $sin^2(\phi) = \frac{1 - \cos(2\phi)}{2}$, kann man $x_3(t)$ auch noch weiter umschreiben:
\begin{align}
    x_3(t) &= \frac{x_{1,0}^2}{2} \cdot (1 - \cos(4 \pi f_1 t)) \\
    &= \frac{x_{1,0}^2}{2} - \frac{x_{1,0}^2}{2} \cdot \cos(4 \pi f_1 t) \nonumber
\end{align}

$\Rightarrow x_3(t) = \underbrace{\frac{x_{1,0}^2}{2}}_{\text{Gleichanteil}} 
- \underbrace{\frac{x_{1,0}^2}{2} \cdot \cos(2\pi \cdot (2f_1) \cdot t)}_{\text{doppelte Frequenz}}$

Somit ist $x_3(t)$ die Summe eines Gleichanteils und eines Signals mit der doppelten Frequenz von $x_1(t)$. 

In Matlab kann $x_3(t)$ wie folgt angezeigt werden:
\begin{sourcecode}
    \caption{Matlab Code zur Berechnung von $x_3(t)$}
\begin{minted}{matlab}
x3 = x1 .* x1;
\end{minted}
\end{sourcecode}

\begin{sourcecode}
    \caption{Matlab Code zum Vergleich von $x_1(t)$ und $x_3(t)$}
\begin{minted}{matlab}
plot(t, x1)
hold on;
plot(t, x3)
hold off;
title('x_1(t) und x_3(t)');
xlabel('Zeit t in [s]');
ylabel('Amplitude');
legend('x_1', 'x_3');
\end{minted}
\end{sourcecode}

\subsection{Lissajous-Figur}
\begin{figure}[H]
    \centering
    \includegraphics[width=0.9\textwidth]{lab1_1_2.png}
    \caption{Lissajous-Figur von $x_1$ und $x_2$}
    \label{fig:lissajous}
\end{figure}
Hier zu sehen ist die Lissajous-Figur von $x_1(t)$ und $x_2(t)$
Um eine Lissajous-Figur zu erstellen, müssen die beiden Signale $x_1(t)$ und $x_2(t)$ um $\varphi = 90^\circ$ phasenverschoben sein, damit die Figur eine geschlossene Form annimmt.
Hierbei wird $x_1(t)$ über $x_2(t)$ geplottet, also $x_1(t)$ auf der x-Achse und $x_2(t)$ auf der y-Achse.
Mit der Lissajous-Figur kann man Phasenverschiebungen zwischen den Signalen erkennen. Früher wurde diese Methode verwendet, um bei Elektronenröhrenoszilloskopen die Phasenverschiebung zwischen zwei Signalen zu bestimmen.
Heutzutage können digitale Oszilloskope die Phasenverschiebung direkt anzeigen, sodass die Lissajous-Figur eher eine anschauliche Darstellung ist.
Zudem kann man die Frequenzverhältnisse der Signale erkennen. 
In diesem Fall ist das Frequenzverhältnis 1:3 (ablesen an der Anzahl der Nullstellen der Signale).
In Matlab wurde die Lissajous-Figur wie folgt dargestellt:
\begin{sourcecode}
    \caption{Matlab Code zur Darstellung der Lissajous-Figur}
\begin{minted}{matlab}
plot(x1,x2)
title('Lissajous-Figur');
xlabel('x_1(t)');
ylabel('x_2(t)');
\end{minted}
\end{sourcecode}

\subsection{Lissajous-Figur mit geringerer Sampling Rate}
\begin{figure}[H]
    \centering
    \includegraphics[width=0.9\textwidth]{lab1_1_3.png}
    \caption{Lissajous-Figur von $x_1$ und $x_2$ mit geringerer Sampling Rate}
    \label{fig:lissajous_wrong}
\end{figure}
Hier zu sehen ist die Lissajous-Figur von $x_1(t)$ und $x_2(t)$, jedoch mit einer geringeren Abtastfrequenz von $f = 500 Hz$.
Da die Abtastfrequenz zu gering ist, kommt es zu Aliasing und das Signal kann nicht mehr korrekt dargestellt werden.
Auf einem Plot von $x_1(t)$ über $x_2(t)$ würden zu niedrige Frequenzen bzw. eine nicht-konstante Frequenz erscheinen.
Das führt zu einer verzerrten Lissajous-Figur wie in Abbildung \ref{fig:lissajous_wrong}.

\begin{sourcecode}
    \caption{Matlab Code zur Signalerzeugung bei zu niedriger Abtastfrequenz}
\begin{minted}{matlab}
t=linspace(0,2,1000);
f1=20e3; 
f2=60e3;

x1 = x1_0 * sin(2*pi*f1*t);
x2 = x2_0 * sin(2*pi*f2*t - pi/2);
\end{minted}
\end{sourcecode}

\begin{sourcecode}[htbp]
    \caption{Matlab Code zur Darstellung der verzerrten Lissajous-Figur}
\begin{minted}{matlab}
plot(x1, x2);
title('Lissajous-Figur mit zu niedrige Abtastfrequenz')
xlabel('x_1(t)');
ylabel('x_2(t)');
\end{minted}
\end{sourcecode}
